\section{Bounded physical state-reduction experiments}
\label{sec:physical-search}

A separate search investigated the primary PA194 table and the new
AX11H and QAX11H tables at the literal transition-table level.
For each table, it tested every unordered pair of controls as an ordinary
merge and a larger family of head-displaced merges. It found no reduction
within these finite-window requirements.

\begin{center}
\begin{tabular}{lrrrr}
\hline
Table & Local windows & Ordinary pairs & Displaced pairs & Hits\\
\hline
Primary PA194 & 10,886 & 18,721 & 32,038 & 0\\
AX11H & 10,909 & 18,721 & 32,038 & 0\\
QAX11H & 11,018 & 18,721 & 32,038 & 0\\
\hline
\end{tabular}
\end{center}

An ordinary merge removes one control, redirects its incoming transitions
and freely resynthesizes the two rows of the retained control. A displaced
merge additionally reverses the movement of every transition into the
removed control, when those incoming movements share a direction. Its
configuration correspondence displaces the represented head by two cells.
Initial controls are included: translation of a blank tape handles a fixed
initial displacement. The local simulation may revisit the retained control.
For each recorded window, it must reproduce the mapped endpoint of one old
transition, including all tape bits, within 100 new transitions while
remaining inside a radius-four window. Exiting that window or exceeding the
local bound rejects a proposal. These restrictions are properties of the
search family, not necessities for every possible correct machine reduction.

The contexts combine symbolic source execution, the interiors of accelerated
constant-symbol scans, physical arithmetic samples, and the exact finite
blank startups. The source checks cover 4,780 outcomes for the primary table
and 4,778 for each hybrid. The startup lengths are 58,154 and 57,922
transitions respectively. Arithmetic interiors were sampled using 3,000
pseudorandom register files and 121 uniform or single-coordinate files,
with 27 block runs per file: 84,267 runs. Sampled values range from zero
through ten, and file lengths range from ten through thirteen.
\emph{This is not a proof that the sampled arithmetic contexts are exhaustive.}
Some symbolic windows also deliberately overapproximate reachable contexts.
Consequently the negative synthesis results are bounded experimental evidence,
not lower bounds on all calibrated machines or on all sound reductions.

The final staged experiment first proposes a locally neutral one-row rewrite
of QAX11H, then attempts every ordinary control-pair merge of the modified
table. The frozen old windows remain in use after the first edit. This is
explicitly a candidate screen; it is not an unproved assertion that those
windows form an invariant of the mutated table. All 358 proposed neutral rewrites were screened. This covered
6,702,118 pair instances, 8,726,731,232 individual action trials,
and 52,929,250 paired action choices, with zero hits. The complete
per-edit counters and all rewrites are supplied in the accompanying results
files.

\subsection{The 193-state runtime splice still fails on both hybrids}

Apply the familiar runtime splice in native numbering: identify control 301
with 88, remove 301, set $(88,0)\mapsto(0,L,51)$, and set
$(51,0)\mapsto(0,R,122)$. Applying this splice to either new hybrid produces
a literal 193-state table. However, both designated blank computations reach
exactly the same configuration after 92 transitions:
\[
 q=35,\qquad h=14,\qquad \operatorname{supp}_1(T)=\{-3,1\}.
\]
The transition at $(35,0)$ writes one, moves right and remains in 35.
Induction on the number of subsequent steps proves an outward infinite ray:
the head always reads the next untouched zero to the right. The supplied
certificate contains every transition of the 92-step prefix and the complete
finite tape support. No timeout is used to infer nonhalting. Thus the H and
AX11 changes do not repair the previously identified initialization failure;
these two 193-state splice attempts are rejected as proof-search machines.
